% Complete editorial constructions; old cited statuses are not altered.
\begin{annexitem}{annex:p02.fbc.finite-syntax}{thm:p02-fbc-finite-syntax}{Finite support, nonvacuous discharge, and finite schemata}
\begin{proof}
A derivation is a finite tree, so it has finitely many leaves and hence
finitely many open assumption formulas. The same tree is permitted
whenever its finite open set is contained in the stated premise set.
This proves monotonicity and finite support without a weakening rule.

For padding, take the given derivation of \(\beta\), add a fresh
assumption leaf \(\alpha\), apply \(\land I\) to obtain
\(\beta\land\alpha\), and apply \(\land E_1\) to conclude \(\beta\).
The new leaf remains an actual open occurrence.
If a derivation of \(\beta\) from \(\Gamma\cup\{\alpha\}\) has
open \(\alpha\)-occurrences, mark all of them with one label fresh
from every existing label, using the source merging operation.
No old discharge uses that fresh label, so those occurrences stay open
until the new root \(\to I\), which discharges them all.
All remaining open formulas belong to \(\Gamma\).
If none occurs, the original derivation is already from \(\Gamma\);
perform the padding just described and discharge its new leaf instead.
In either case the discharge is nonvacuous.
For the reverse implication, take a fresh-labelled copy of the
derivation of \(\alpha\to\beta\), add an assumption \(\alpha\),
and apply \(\to E\). The resulting open set is contained in
\(\Gamma\cup\{\alpha\}\).

Uniform formula substitution fixes \(\bottom\) and commutes with
each primitive constructor. Apply it to the formula at every node
of a finite derivation, leaving matching assumption/rule labels intact.
Each assumption is still an assumption; each \(\land I,\land E\),
\(\lor I,\lor E,\to I,\to E,\neg I,\neg E,\bottom E\) and RAA
instance is the same rule with its formula metavariables substituted.
Every discharged labelled occurrence is still present and has the
substituted bracketed formula. Thus no discharge becomes vacuous,
even if different old formulas obtain equal substitutions.
Substitution does not act again inside an inserted formula.
When copies are assembled, finitely many labels can be renamed to
distinct natural numbers greater than every receiving label.
This is a finite construction, not a choice of names across an
arbitrary family.

Now take a tautological schema \(\theta(p_0,\ldots,p_{k-1})\) in
the fixed source alphabet. Its truth value depends only on these
finitely many atoms, by structural induction. Validity for their
binary assignments therefore makes it valid for all assignments on
the source countable alphabet. The already proved
\cref{thm:prop-completeness}, with empty premise set, gives a closed
finite derivation of \(\theta\). That theorem's actual proof uses
the source countable Lindenbaum enumeration and truth lemma,
not Boolean quotients or arbitrary-language completeness.
That enumeration requires no Choice: encode \(p_n\) by its
natural-number binary digits with fixed delimiters, encode the
finitely many other tokens by fixed characters, and enumerate the
resulting finite words by length and then lexicographically.
Retaining the well-formed encodings lists every source formula.
This construction concerns the given indexed source atoms, not \(P\).
Assign prescribed \(P\)-formulas to the displayed \(k\) atoms and
\(\bottom\) to every unused source atom. This explicitly defines
one total substitution even when \(P=\varnothing\).
The preceding rule induction gives the desired closed instance.
Only this one proof of the selected schema is used; no simultaneous
selection of proofs for all schemata is stipulated.
\end{proof}
\end{annexitem}

\begin{annexitem}{annex:p02.fbc.boolean-soundness}{thm:p02-fbc-boolean-soundness}{Boolean-valued soundness for every ND rule}
\begin{proof}
The complement/order calculus of \cref{thm:p01-bpi-duality} follows
from Boolean axioms and uses no representation or choice.
In particular
\[
 r\wedge s=0_B\quad\Longleftrightarrow\quad r\leq_B\neg s.
\]
We will also use the residuation identity
\[
 x\wedge a\leq_B b
       \quad\Longleftrightarrow\quad x\leq_B\neg a\vee b.
\]
For the forward direction distribute
\(x=x\wedge(a\vee\neg a)\):
the first summand is below \(b\), the second below \(\neg a\).
For the reverse direction meet \(x\leq_B\neg a\vee b\) with \(a\);
the upper bound becomes \(a\wedge b\leq_B b\).
These calculations remain valid when \(0_B=1_B\).

Induct on the derivation, using the meet of its actual finite open
assumption values, not of the possibly infinite permitted set.
An assumption leaf has value \(a\) and open meet \(a\).
For \(\land I\), if the premise meets are \(x,y\), then
\(x\leq a\), \(y\leq b\) imply \(x\wedge y\leq a\wedge b\).
Overlapping open sets cause no problem, since repeats are idempotent.
For \(\land E_1,\land E_2\), the premise bound
\(x\leq a\wedge b\) implies the respective bound to \(a,b\).
For \(\lor I_1,\lor I_2\), use respectively
\(a\leq a\vee b\) and \(b\leq a\vee b\).

For \(\to E\), the combined open meet is below both
\(\neg a\vee b\) and \(a\), hence below
\(a\wedge(\neg a\vee b)=a\wedge b\leq b\).
For \(\neg E\), it is below \(a\wedge\neg a=0_B\).
For \(\bottom E\), a bound below \(0_B\) is below every element.

For \(\to I\), let \(x\) be the meet of remaining open
assumptions and \(a\) the value of the discharged formula.
The premise open meet is \(x\wedge a\), even if some
undischarged occurrences of the same formula remain in \(x\).
The induction hypothesis gives \(x\wedge a\leq b\);
residuation gives \(x\leq\neg a\vee b\).
For \(\neg I\), it gives \(x\wedge a\leq0_B\), hence
\(x\leq\neg a\).
For RAA the discharged formula has value \(\neg a\);
the premise gives \(x\wedge\neg a\leq0_B\), hence
\(x\leq\neg\neg a=a\).
This is exactly where Boolean, rather than merely intuitionistic,
complementation supports the classical rule.

For \(\lor E\), let the major premise meet be \(x\) and let
the subsidiary remaining-open meets be \(y,z\).
With \(a,b\) the disjunct values and \(c\) the conclusion value,
the induction hypotheses are
\[
 x\leq a\vee b,\qquad y\wedge a\leq c,\qquad
 z\wedge b\leq c.
\]
The final open meet \(x\wedge y\wedge z\) is below
\[
 (a\vee b)\wedge y\wedge z
       =(a\wedge y\wedge z)\vee(b\wedge y\wedge z)
       \leq c.
\]
Both labelled branch assumptions are accounted for before removal.
These cases exhaust the actual rule sheet.

A closed tree has empty open meet \(1_B\), so its conclusion value
is \(1_B\). If that conclusion is \(\alpha\leftrightarrow\beta\),
both \(\neg a\vee b\) and \(\neg b\vee a\) are one.
Residuation with \(x=1_B\) gives \(a\leq b\) and \(b\leq a\),
hence \(a=b\). No completeness assertion was used.
\end{proof}
\end{annexitem}

\begin{annexitem}{annex:p02.fbc.quotient}{thm:p02-fbc-quotient}{Congruence and all Boolean identities of the syntactic quotient}
\begin{proof}
The value of a biconditional in \(\mathbf2\) is one exactly when
its two arguments have equal values. Thus the following finite schemata
are tautologies: \(\alpha\leftrightarrow\alpha\);
\((\alpha\leftrightarrow\beta)\to(\beta\leftrightarrow\alpha)\);
and
\[
 ((\alpha\leftrightarrow\beta)\land
   (\beta\leftrightarrow\gamma))
       \to(\alpha\leftrightarrow\gamma).
\]
The first is equality with itself, the second reverses equality,
and the third is transitivity of equal binary values.
\Cref{thm:p02-fbc-finite-syntax} gives closed ND instances.
Conjunction introduction and implication elimination on the given
closed equivalences establish reflexivity, symmetry and transitivity
of \(\sim_P\).

For congruence, the finite schema
\[
 (\alpha\leftrightarrow\alpha')
                 \to(\neg\alpha\leftrightarrow\neg\alpha')
\]
is tautological, since complementation is a function on binary values.
For each \(\circ\in\{\land,\lor,\to\}\), so is
\[
 ((\alpha\leftrightarrow\alpha')\land
   (\beta\leftrightarrow\beta'))
       \to((\alpha\circ\beta)\leftrightarrow
                              (\alpha'\circ\beta')).
\]
Equal pairs of inputs have equal outputs under each truth function.
Closed instances and \(\land I,\to E\) give all congruence clauses.

For explicitness, all the required Boolean identities are obtained
as biconditionals of the following formula pairs:
\[
\begin{array}{c}
 \alpha\land\beta,\ \beta\land\alpha\\
 \alpha\lor\beta,\ \beta\lor\alpha\\
 (\alpha\land\beta)\land\gamma,\ \alpha\land(\beta\land\gamma)\\
 (\alpha\lor\beta)\lor\gamma,\ \alpha\lor(\beta\lor\gamma)\\
 \alpha\land\alpha,\ \alpha\qquad
       \alpha\lor\alpha,\ \alpha\\
 \alpha\land(\alpha\lor\beta),\ \alpha\\
 \alpha\lor(\alpha\land\beta),\ \alpha\\
 \alpha\land\topv,\ \alpha\qquad\alpha\lor\bottom,\ \alpha\\
 \alpha\land\bottom,\ \bottom\qquad\alpha\lor\topv,\ \topv\\
 \alpha\land(\beta\lor\gamma),\
       (\alpha\land\beta)\lor(\alpha\land\gamma)\\
 \alpha\lor(\beta\land\gamma),\
       (\alpha\lor\beta)\land(\alpha\lor\gamma)\\
 \alpha\land\neg\alpha,\ \bottom\qquad
       \alpha\lor\neg\alpha,\ \topv\\
 \neg\neg\alpha,\ \alpha\\
 \alpha\to\beta,\ \neg\alpha\lor\beta\\
 \neg(\alpha\land\beta),\ \neg\alpha\lor\neg\beta\\
 \neg(\alpha\lor\beta),\ \neg\alpha\land\neg\beta .
\end{array}
\]
Here commas separate the two members of a pair, not premises.
We verify that every pair has equal binary values before invoking
finite schematic derivability.
In \(\mathbf2\), conjunction is minimum and disjunction is maximum.
Commutativity follows because exchanging the two inputs does not
change their lower or upper value. Both associations give the least
or greatest of the same three inputs. Idempotence uses one input twice.
For absorption, \(u\leq\max(u,v)\) and \(\min(u,v)\leq u\);
the indicated outer minimum or maximum is consequently \(u\).
The four bound laws follow from \(0\leq u\leq1\).
For meet distributing over join, \(u=0\) makes both sides zero;
\(u=1\) makes both sides \(\max(v,w)\).
For the dual distribution, \(u=1\) makes both sides one;
\(u=0\) makes both sides \(\min(v,w)\).
For \(u=0\), the complement is one; for \(u=1\), it is zero.
In both cases their minimum is zero and their maximum is one.
Double negation is \(1-(1-u)=u\).
For the complement of a meet, the two sides have value one when
both inputs are zero, zero when both are one, and one when they differ.
For the complement of a join, those three cases give respectively
one, zero and zero on both sides. These exhaust the four assignments.
The implication pair is its specified binary truth table:
when \(u=0\) both sides are one, and when \(u=1\) both equal \(v\).
Each biconditional is therefore tautological, and
\cref{thm:p02-fbc-finite-syntax} gives its closed proof.

Replacement forms the set of all equivalence classes.
Congruence makes each displayed operation a well-defined function:
any two representatives of one input have equivalent outputs, and
the same holds for the two inputs of a binary operation.
Representatives of the finitely many classes in a particular identity
exist by finite witness selection in ZF; no selector on all classes
is required. The closed biconditionals just proved show that each
listed identity holds in the quotient. They include every bounded
distributive lattice and complement axiom, making it a Boolean algebra.
The last implication pair gives the asserted implication interpretation.

Finally a total valuation on \(P\) exists without Choice: assign
zero to every atom, using the empty function when \(P\) is empty.
Its recursive extension sends \(\bottom\) to zero and
\(\topv=\neg\bottom\) to one. If their classes were equal,
Boolean-valued soundness in \(\mathbf2\) would make these values equal,
which is false. Thus \(L_P\) is nontrivial in every case.
\end{proof}
\end{annexitem}

\begin{annexitem}{annex:p02.fbc.free}{thm:p02-fbc-free}{The universal map and syntactic order}
\begin{proof}
Fix a Boolean algebra \(B\) and an assignment \(f:P\to B\).
Structural recursion on finite formulas defines
\(\llbracket\alpha\rrbracket_f\).
If \(\alpha\sim_P\beta\), the given closed proof of their
biconditional and \cref{thm:p02-fbc-boolean-soundness} give equal
values. The relation
\[
 h_f([\alpha]_P)=\llbracket\alpha\rrbracket_f
\]
therefore specifies exactly one value at every class, independently
of all its representatives. Replacement forms its graph as a function;
it does not choose a representative for each class.
The evaluation clauses preserve zero, one, complement, meet and join,
and on an atomic class give \(h_f([p]_P)=f(p)\).
Thus \(h_f\) is a homomorphism with \(h_f\eta_P=f\).

If \(h\) is any such homomorphism, induction on \(\alpha\) gives
\(h([\alpha]_P)=\llbracket\alpha\rrbracket_f\).
The atomic case is \(h\eta_P=f\); the bottom case is preservation
of zero. Negation, meet and join use their quotient definitions and
homomorphism laws. Implication uses the already established identity
\([\alpha\to\beta]=\neg[\alpha]\vee[\beta]\).
These exhaust the grammar, proving uniqueness.
No nontriviality of \(B\) was required. If it is trivial all values
and all the preservation equations still exist and agree.

For distinct atoms \(p,q\), the assignment to \(\mathbf2\) taking
value one precisely at \(p\) and zero at every other atom is a set
function defined by Separation. Its induced homomorphism sends
\([p]\) to one and \([q]\) to zero, so those classes differ.
Hence \(\eta_P\) is injective.
If \(P=\varnothing\), induction on a formula shows that its class
is either \(0_L\) or \(1_L\): bottom starts at zero; complement
interchanges zero and one; meet and join of these two stay among them;
and implication is complement followed by join.
Their distinctness was proved in the quotient item.
The same existence/uniqueness proof applies to the sole empty assignment
into any target algebra.

For the order assertion, the two finite schemata
\[
 ((\alpha\land\beta)\leftrightarrow\alpha)\to(\alpha\to\beta),
 \qquad
 (\alpha\to\beta)\to((\alpha\land\beta)\leftrightarrow\alpha)
\]
are tautological. Indeed \(\min(u,v)=u\) iff \(u=0\) or \(v=1\),
and that is exactly the condition for \(u\to v=1\).
Their closed instances exist by finite schematic derivability.
The equation \([\alpha]\wedge[\beta]=[\alpha]\) is, by the
quotient definition, the existence of the first inner biconditional
as a closed proof. Modus ponens with the first schema gives
\(\varnothing\vdash_P\alpha\to\beta\).
Conversely that closed implication and the second schema give
the biconditional and hence the defining order equation.
\end{proof}
\end{annexitem}

\begin{annexitem}{annex:p02.fbc.filter}{thm:p02-fbc-filter}{The exact premise-filter correspondence}
\begin{proof}
The empty sequence shows \(1_L\in F_\Gamma\), and a one-member
sequence shows \([\gamma]\in F_\Gamma\) for each \(\gamma\in\Gamma\).
Transitivity of order proves upward closure.
Concatenating two witnessing finite sequences gives a meet below the
meet of their respective upper bounds, so proves meet closure.
Every filter containing the generators contains each finite meet and
everything above it. Thus \(F_\Gamma\) is exactly their least filter,
also the generated filter of \cref{thm:p01-bpi-duality}.

For the forward derivability-to-membership implication, interpret the
calculus in \(L_P\) by the assignment \(p\mapsto[p]_P\).
Induction on a formula gives
\(\llbracket\alpha\rrbracket_\eta=[\alpha]_P\): the bottom,
complement, meet and join cases are definitions, and the implication
case uses its proved quotient identity.
If \(D\) derives \(\varphi\) with open formulas in \(\Gamma\),
Boolean-valued soundness gives
\[
 \bigwedge_{\delta\in\operatorname{oa}(D)}[\delta]_P
                                  \leq_L[\varphi]_P .
\]
Its finite open set can be listed in a finite sequence in ZF;
this one sequence witnesses membership.

Conversely take a witnessing finite sequence
\(\gamma_1,\ldots,\gamma_n\), and let \(\tau\) be its
right-associated formula conjunction.
Its class is the witnessing Boolean meet, including
\([\tau]=[\topv]=1_L\) when \(n=0\).
The inequality to \([\varphi]\) and the syntactic order theorem
give a closed proof of \(\tau\to\varphi\).
For \(n>0\), fresh assumption leaves for the sequence entries and
successive \(\land I\) from the right build \(\tau\).
If entries repeat, distinct labelled leaves are allowed; their
permitted formula set is still contained in \(\Gamma\).
For \(n=0\), assume \(\bottom\) as one fresh leaf and discharge
that very leaf by \(\neg I\); this gives the closed proof
of \(\neg\bottom=\topv=\tau\).
Freshen the labels of the closed implication proof and apply
\(\to E\) to these two derivations.
No extra assumption is introduced; the resulting open set is
contained in the finite sequence's members, and hence in \(\Gamma\).
This proves the reverse correspondence with actual ND trees.

Taking \(\varphi=\bottom\) shows that inconsistency is equivalent
to \(0_L\in F_\Gamma\). A filter containing zero contains every
element by upward closure; conversely the whole algebra contains zero.
Thus properness is exactly consistency.
When \(\Gamma=\varnothing\), only the empty sequence is permitted,
so its upper bounds are exactly \(1_L\), proving
\(F_\varnothing=\{1_L\}\). Nontriviality makes this filter proper.
\end{proof}
\end{annexitem}

\begin{annexitem}{annex:p02.fbc.separation}{thm:p02-fbc-separation}{The avoiding filter and every truth-lemma constructor}
\begin{proof}
Write \(F=F_\Gamma\) and \(a=[\varphi]_P\).
The correspondence gives \(a\notin F\); this also implies \(F\)
proper, since an improper filter contains everything.
The displayed \(H\) contains one by witness \(f=1_L\), and contains
\(\neg a\) by the same witness. Every \(f\in F\) belongs to \(H\),
since \(f\wedge\neg a\leq f\).
It is upward closed. If \(b,c\) have witnesses \(f,g\in F\),
then \(f\wedge g\in F\) and
\[
 (f\wedge g)\wedge\neg a
        =(f\wedge\neg a)\wedge(g\wedge\neg a)\leq b\wedge c .
\]
Thus \(H\) is a filter.
If it contained zero, some \(f\in F\) would satisfy
\(f\wedge\neg a=0_L\), hence \(f\leq a\) by Boolean
complement/order calculus. Upward closure would put \(a\) in \(F\),
contrary to the premise. This proves properness before extension.

Assume BPI and use \cref{thm:p01-bpi-equivalence} once to extend
this prescribed proper filter to an ultrafilter \(U\).
Since \(\neg a\in U\), properness forbids \(a\in U\).
The ZF duality theorem \cref{thm:p01-bpi-duality} proves that
\(\chi_U:L_P\to\mathbf2\) preserves both constants and all
Boolean operations. Define \(v_U=\chi_U\eta_P\).
This is one specified function on \(P\), not separate choices
of a truth value for each atom.

Induct on the formula for the truth lemma.
For an atom it is the definition.
For bottom, both sides are zero by preservation of \(0_L\).
For negation, both sides are the Boolean complement of the value
of the immediate subformula, by the induction hypothesis and
preservation of complement.
For conjunction and disjunction, use the corresponding two
induction hypotheses and preservation of meet and join.
For implication, use
\[
 \chi_U([\alpha\to\beta])
     =\chi_U(\neg[\alpha]\vee[\beta])
     =\neg\chi_U([\alpha])\vee\chi_U([\beta]),
\]
which is precisely its recursive binary truth value.
The defined \(\topv\) case follows from negation of bottom.
These are all primitive constructors and no infinitary clause
is introduced.
For \(\gamma\in\Gamma\), its class belongs to \(F\subseteq U\)
and so its truth value is one. The class \(a\) is excluded,
so \(\varphi\) has value zero.
The proof includes empty \(\Gamma\) and empty \(P\);
if the latter occurs, \(v_U\) is the unique empty assignment.
\end{proof}
\end{annexitem}

\begin{annexitem}{annex:p02.fbc.completeness}{thm:p02-fbc-completeness}{Completeness, finite support and equivalence to BPI}
\begin{proof}
For soundness take \(B=\mathbf2\) in Boolean-valued soundness.
A valuation satisfying all permitted premises gives value one
to the finite open meet, hence to the conclusion.
This works in ZF alone.

Assume BPI. If \(\Gamma\not\vdash_P\varphi\), the preceding
construction gives a total valuation satisfying \(\Gamma\) and
falsifying \(\varphi\). This is exactly
\(\Gamma\not\models_P\varphi\); contraposition gives completeness.
When \(\Gamma\models_P\varphi\), its obtained derivation has
finite open set \(\Gamma_0\subseteq\Gamma\).
The same derivation is from \(\Gamma_0\), and soundness gives
\(\Gamma_0\models_P\varphi\).

If every finite subset of \(\Gamma\) is satisfiable but
\(\Gamma\) has no satisfying valuation, then
\(\Gamma\models_P\bottom\).
Completeness gives a finite derivation of bottom, and its finite
open set \(\Gamma_0\) cannot be satisfiable by soundness.
This contradicts the premise. The converse follows by using any
model of \(\Gamma\) for each subset.

For the choice-strength converse, assume over ZF the completeness
scheme for every set \(P\), without initially assuming BPI.
The same bottom/finite-support argument, now using that scheme
and the already choice-free soundness lemma, proves compactness
for every such \(P\).
Restrict to the \(\neg,\land,\lor\) fragment and interpret the
separate top constant of
\cref{def:p01-bpi-propositional-compactness} by the source
abbreviation \(\neg\bottom\).
Induction on formulas shows that this translation preserves every
truth value. It also covers the empty atom set, since no atom
is needed to form either constant.
More explicitly, translate a finitely satisfiable set \(\Sigma\)
in the P1 language. For a finite subset of its translated image,
fix one preimage per member; this is finite witness selection in ZF.
Those finitely many preimages have a common valuation, which satisfies
the selected translated formulas by truth preservation.
The image is therefore finitely satisfiable. The compactness just
obtained supplies a valuation of the whole image; truth preservation
then makes it a valuation of every member of \(\Sigma\).
Thus we obtain exactly the arbitrary-set-variable compactness
principle, not merely its countable restriction.
The complete converse construction of
\cref{thm:p01-bpi-propositional-compactness} then gives BPI.
Together with the proved BPI-to-completeness direction this
establishes the scheme equivalence over ZF.

No unrestricted completeness was used in the quotient or free
property proof. The countable source theorem was used only for
closed finite schemata before those constructions, and its
independent enumeration proof remains within its fixed alphabet.
Nothing here makes a general algorithm or well-order for an
arbitrary set of atoms; the precise algorithmic boundary is proved in
\cref{thm:p02-w02-p02-017-no-coding,thm:p02-w02-p02-017-bpi-boundary}.
\end{proof}
\end{annexitem}
